%======================================================================
\chapter{The main theorems}\label{ch:main}
%======================================================================

This chapter assembles Parts~I--III and proves the theorems announced in the introduction:
Theorems~\ref{thm:intro-unique} and~\ref{thm:intro-homog} and Corollary~\ref{cor:intro}. All the
work has been done; what remains is to put the pieces together, and to be careful about the
ambient spaces. Throughout, $X$ and $X'$ are metric spaces.

Recall (\vone{Definition~2.14}) that a \emph{pseudo-arc} is a nondegenerate, snake-like,
hereditarily indecomposable continuum $M\subseteq X$.

\formal{\lean{HPC.Plane.IsPseudoArc}: the conjunction of \lean{HPC.Plane.IsSnakeLike} (which
contains ``$M$ is a continuum''), \lean{HPC.Plane.IsHereditarilyIndecomposable} and
\lean{Set.Nontrivial}.}

\section{Uniqueness}\label{sec:main-unique}

\begin{theorem}[Bing 1951, Theorem~1, the ``in fact'' clause]\label{thm:unique-points}
Let $M\subseteq X$ and $M'\subseteq X'$ be pseudo-arcs. Let $p,q\in M$ with $q\notin\cmp Mp$, and
let $p',q'\in M'$ with $q'\notin\cmp{M'}{p'}$. Then there is a homeomorphism $h\colon M\to M'$
with $h(p)=p'$ and $h(q)=q'$.
\end{theorem}

\begin{proof}
Regard $M$ and $M'$ as compact metric spaces (with the metrics of $X$ and $X'$).

\emph{Step 1: descriptions.} By Theorem~\ref{thm:desc-exists}(a) there is a description
$C^{(0)},C^{(1)},\dots$ from $p$ to $q$ \emph{in the space $M$} whose described set is the whole
space $M$, and likewise a description from $p'$ to $q'$ in the space $M'$ whose described set is
all of $M'$.

\emph{Step 2: the end points are distinct.} A pseudo-arc has at least two points, so $p\ne q$
and $p'\ne q'$ by Lemma~\ref{lem:ne-of-composant}.

\emph{Step 3: two descriptions.} The spaces $M$ and $M'$ are compact, and the two descriptions go
between distinct points. By Theorem~\ref{thm:two-desc} there is a homeomorphism $T$ from the
described set of the first description onto the described set of the second, with $T(p)=p'$ and
$T(q)=q'$. These described sets are $M$ and $M'$, so $h=T$ is the required homeomorphism.
\end{proof}

\formal{\lean{HPC.Plane.theorem1_HIC}. Formally the described sets are subsets
$\Delta_1.M\subseteq M$ and $\Delta_2.M\subseteq M'$ of the spaces $M$, $M'$, equal to the whole
spaces; the homeomorphism of \lean{HPC.theorem12} is composed with the identifications
$M\cong\Delta_1.M$ and $\Delta_2.M\cong M'$.}

\begin{remark}
Note how little of Theorem~\ref{thm:desc-exists} is used: only part~(a), the description
inside $M$ itself. Working inside $M$ has the advantage that the ambient space is compact, which
is exactly what Theorem~\ref{thm:two-desc} asks for. Part~(b) would give descriptions in $X$ and
$X'$; to use those one would need their links to lie in compact sets
(Corollary~\ref{cor:two-desc-K}), which is true for the lifted descriptions of
Section~\ref{sec:lift} only when $X$ is, for example, a proper space.
\end{remark}

\begin{theorem}[Uniqueness of the pseudo-arc]\label{thm:unique}
Any two pseudo-arcs $M\subseteq X$ and $M'\subseteq X'$ are homeomorphic.
\end{theorem}

\begin{proof}
$M$ is nondegenerate, so it has a point $p$; by Corollary~\ref{cor:other-composant} there is
$q\in M$ with $q\notin\cmp Mp$. Likewise choose $p'\in M'$ and $q'\in M'$ with
$q'\notin\cmp{M'}{p'}$. Apply Theorem~\ref{thm:unique-points}.
\end{proof}

\formal{\lean{HPC.Plane.theorem1_HIC'}; the same statement with the three defining properties
written out is \lean{HPC.Plane.theorem1_snake} (Bing, \emph{Snake-like continua}, Theorem~1).}

Theorem~\ref{thm:intro-unique} of the introduction is Theorems~\ref{thm:unique-points}
and~\ref{thm:unique} together.

\section{Homogeneity}\label{sec:main-homog}

\begin{theorem}[Homogeneity of the pseudo-arc]\label{thm:homog}
Let $M\subseteq X$ be a pseudo-arc and $p,q\in M$. There is a homeomorphism $h\colon M\to M$ with
$h(p)=q$.
\end{theorem}

\begin{proof}
By Corollary~\ref{cor:other-composant} choose $r\in M$ with $r\notin\cmp Mp$ and $s\in M$ with
$s\notin\cmp Mq$. Apply Theorem~\ref{thm:unique-points} with $M'=M$, to the pairs $(p,r)$ and
$(q,s)$: there is a homeomorphism $h\colon M\to M$ with $h(p)=q$ (and $h(r)=s$).
\end{proof}

\formal{\lean{HPC.Plane.homogeneous_of_isPseudoArc}.}

This is Theorem~\ref{thm:intro-homog}. The proof shows more: for any two pairs $(p,r)$ and
$(q,s)$ of points in different composants, some homeomorphism of $M$ onto itself takes $p$ to $q$
and $r$ to $s$ at the same time.

\begin{remark}
It is instructive to see why the argument does not prove that the arc $[0,1]$ is homogeneous,
which is false. The arc is snake-like but decomposable, so Chapter~\ref{ch:descriptions} does not
apply: fine chains covering $[0,1]$ need not be crooked in coarse ones
(Theorem~\ref{thm:eventually-crooked} fails). And even the input of Chapter~\ref{ch:composants}
fails: for $0<p<1$ the composant of $p$ is all of $[0,1]$, so there is no partner point $r$.
\end{remark}

\section{Consequences}\label{sec:main-cons}

\begin{corollary}\label{cor:plane-unique}
\begin{enumerate}[label=(\alph*)]
\item There is a pseudo-arc $M_0$ in the plane $\R^2$.
\item Every pseudo-arc, in any metric space, is homeomorphic to $M_0$.
\end{enumerate}
\end{corollary}

\begin{proof}
(a) is \vone{Theorem~5.18}: the plane description of \vone{Theorem~5.17} goes between two distinct
points and all its links lie in a compact rectangle, so its described set is a pseudo-arc by
\vone{Theorem~3.8}. (b) follows from Theorem~\ref{thm:unique}.
\end{proof}

\formal{\lean{HPC.Plane.isPseudoArc_plane} and \lean{HPC.Plane.pseudoArc_unique}.}

So ``the pseudo-arc'' is well defined up to homeomorphism, and it can be drawn in the plane. This
is Corollary~\ref{cor:intro}(a).

\begin{lemma}\label{lem:sub-pseudoarc}
Every nondegenerate subcontinuum $K$ of a pseudo-arc $M$ is a pseudo-arc.
\end{lemma}

\begin{proof}
$K$ is a nondegenerate continuum by assumption. Every subcontinuum of $K$ is a subcontinuum of
$M$, hence indecomposable; so $K$ is hereditarily indecomposable. Given $\varepsilon>0$, an
$\varepsilon$-chain covering $M$ also covers $K$; so $K$ is snake-like.
\end{proof}

\formal{\lean{HPC.Plane.isPseudoArc_of_subcontinuum}.}

\begin{corollary}\label{cor:sub-homeo}
Every nondegenerate subcontinuum of a pseudo-arc $M$ is homeomorphic to $M$.
\end{corollary}

\begin{proof}
Lemma~\ref{lem:sub-pseudoarc} and Theorem~\ref{thm:unique}.
\end{proof}

\formal{\lean{HPC.Plane.homeomorphic_subcontinuum}.}

This is Corollary~\ref{cor:intro}(b). Together with Corollary~\ref{cor:plane-unique} it shows
that the pseudo-arc is homeomorphic to each of its nondegenerate subcontinua, although it
contains no arc (an arc is decomposable).

\section{End points}\label{sec:main-endpoints}

Bing's paper \emph{Snake-like continua} characterizes the pseudo-arc among snake-like continua by
a property of its points. We prove this characterization, since it follows quickly from what we
have.

\begin{definition}\label{def:endpoint}
Let $M\subseteq X$ be a continuum and $p\in M$.
\begin{enumerate}[label=(\alph*)]
\item $p$ is an \emph{end point} of $M$ if for every $\varepsilon>0$ there is an $\varepsilon$-chain $E$ of
length $n\ge1$ covering $M$ such that $p$ lies in the first link $E_0$ and in no other link.
\item Points $p,q\in M$ are \emph{opposite end points} of $M$ if for every $\varepsilon>0$ there is an
$\varepsilon$-chain covering $M$ which goes from $p$ to $q$.
\end{enumerate}
\end{definition}

\formal{\lean{HPC.Plane.IsEndPoint}, \lean{HPC.Plane.OppositeEndPoints} (in
\lean{HPC/Plane/EndPoints51.lean}).}

\begin{example}
The points $0$ and $1$ are opposite end points of the arc $[0,1]$: cover it by the chain of
intervals $(-\delta,2\delta),(\delta,3\delta),\dots$ of length $3\delta$ and step $\delta$, where $\delta=1/m$
is small, cut down to the links meeting $[0,1]$. A point $p$ with $0<p<1$ is not an end point of
$[0,1]$; this follows from Lemma~\ref{lem:endpoint-nested} below, since $[0,p]$ and $[p,1]$ are
subcontinua containing $p$, neither of which contains the other.
\end{example}

\begin{proposition}\label{prop:opposite}
Let $M$ be a pseudo-arc and $p,q\in M$ with $q\notin\cmp Mp$. Then $p$ and $q$ are opposite end
points of $M$. In particular every point of a pseudo-arc is an end point of it.
\end{proposition}

\begin{proof}
By Theorem~\ref{thm:desc-exists}(b) there is a description $D^{(0)},D^{(1)},\dots$ in $X$ from $p$
to $q$ with described set $M$. Given $\varepsilon>0$, choose $i$ with $\varepsilon_i<\varepsilon$. The
chain $D^{(i)}$ is an $\varepsilon$-chain by~\vone{(D4)}, it covers $M$ because $M$ is the described
set, and it goes from $p$ to $q$ by~\vone{(D1)}.

For the second statement let $p\in M$; by Corollary~\ref{cor:other-composant} there is $q\in M$
with $q\notin\cmp Mp$, and the chains just found have $p$ in the first link only.
\end{proof}

\formal{\lean{HPC.Plane.oppositeEndPoints_of_composants}.}

Here part~(b) of Theorem~\ref{thm:desc-exists} is really used: the chains must be chains of $X$.

\begin{lemma}\label{lem:endpoint-nested}
Let $p$ be an end point of a continuum $M$, and let $A,B\subseteq M$ be continua containing $p$.
Then $A\subseteq B$ or $B\subseteq A$.
\end{lemma}

\begin{proof}
Suppose not: there are $a\in A\setminus B$ and $b\in B\setminus A$. Since $A$ and $B$ are compact,
$\delta=\min(\dist(a,B),\dist(b,A))>0$. Let $E$ be a $\delta$-chain of length $n$ covering $M$ with
$p$ in $E_0$ only.

The links of $E$ that meet $A$ form an interval of indices $[0,i_A]$: it is an interval by
Lemma~\ref{lem:cover-meeting}(a) (applied to the connected set $A$, which is covered by $E$), and
it contains $0$ because $p\in A\cap E_0$. Likewise the links meeting $B$ form an interval
$[0,i_B]$. Suppose $i_A\le i_B$ (the other case is symmetric, exchanging the roles of $a$ and $b$).
The point $a$ lies in some link $E_j$; this link meets $A$, so $j\le i_A\le i_B$, and therefore $E_j$
also meets $B$, say at $z$. Then $\dist(a,B)\le d(a,z)\le\diam E_j<\delta$, a contradiction.
\end{proof}

\formal{\lean{HPC.Plane.condB_of_endPoint}.}

\begin{lemma}\label{lem:not-HI-not-end}
If a continuum $M$ is not hereditarily indecomposable, then some point of $M$ is not an end point
of $M$.
\end{lemma}

\begin{proof}
Let $K\subseteq M$ be a decomposable subcontinuum, $K=A\cup B$ with $A,B$ subcontinua of $K$
different from $K$. Since $K$ is connected and $A,B$ are nonempty closed sets covering it, they
meet; let $p\in A\cap B$. If $p$ were an end point of $M$, Lemma~\ref{lem:endpoint-nested} would
give, say, $A\subseteq B$, and then $K=A\cup B=B$, a contradiction.
\end{proof}

\formal{\lean{HPC.Plane.exists_not_endPoint_of_not_HI}.}

\begin{theorem}[Bing 1951, \emph{Snake-like continua}, Theorem~16]\label{thm:endpoints}
A nondegenerate snake-like continuum $M$ is a pseudo-arc if and only if every point of $M$ is an
end point of $M$.
\end{theorem}

\begin{proof}
If $M$ is a pseudo-arc, every point is an end point by Proposition~\ref{prop:opposite}.
Conversely, if every point is an end point, $M$ is hereditarily indecomposable by
Lemma~\ref{lem:not-HI-not-end}, and hence a pseudo-arc.
\end{proof}

\formal{\lean{HPC.Plane.isPseudoArc_iff_endPoints}.}

\section{What the classification uses}\label{sec:main-uses}

The later volumes use the results of this volume in the following forms.
\begin{itemize}
\item \emph{Uniqueness} (Theorem~\ref{thm:unique}, Corollary~\ref{cor:plane-unique}): to recognize
a continuum as ``the'' pseudo-arc it suffices to check the three defining properties. This is how
continua arising later (for instance the indecomposable homogeneous plane continua) are
identified with the pseudo-arc: one proves that they are snake-like and hereditarily
indecomposable.
\item \emph{Homogeneity} (Theorem~\ref{thm:homog}): the pseudo-arc appears in the list of
homogeneous plane continua, so it must be shown to be homogeneous.
\item \emph{Subcontinua} (Lemma~\ref{lem:sub-pseudoarc}, Corollary~\ref{cor:sub-homeo}).
\end{itemize}

\begin{remark}[On the hypotheses]
Two statements in the literature need an extra hypothesis, which the formal proof adds.
\begin{enumerate}[label=(\alph*)]
\item Bing's Theorem~12 for descriptions whose described sets are merely compact is false
(Remark after Corollary~\ref{cor:two-desc-K}); one needs the links to lie in a compact set.
\item In the same way, ``the described set of a description from $P$ to $Q\ne P$ is a pseudo-arc''
needs the links to lie in a compact set (\vone{Theorem~3.8} and \vone{Remark~3.5}); the formal
record of the counterexample is \lean{HPC.Plane.isPseudoArc_of_description_false}.
\end{enumerate}
Neither affects the theorems of this chapter, whose hypotheses only involve pseudo-arcs, which
are compact by definition.
\end{remark}

\section*{Exercises}
\addcontentsline{toc}{section}{Exercises}

\begin{exercise}\label{ex:arc-endpoints}
Show directly from Definition~\ref{def:endpoint} that the end points of $[0,1]$ are exactly $0$
and $1$. (Hint: use Lemma~\ref{lem:endpoint-nested} with $A=[0,p]$ and $B=[p,1]$.)
\end{exercise}

\begin{exercise}
Let $p$ be an end point of a continuum $M$ and let $K\subseteq M$ be a continuum containing $p$. Show
that $p$ is an end point of $K$. (Hint: Lemma~\ref{lem:cover-meeting}(a).) Deduce again that every
point of a pseudo-arc is an end point of each subcontinuum containing it.
\end{exercise}

\begin{exercise}
Let $M$ be a pseudo-arc and $p\in M$. Show that the homeomorphisms $h\colon M\to M$ with $h(p)=p$
act transitively on $M\setminus\cmp Mp$. (Hint: Theorem~\ref{thm:unique-points} with
$p'=p$.)
\end{exercise}

\begin{exercise}
Show that a pseudo-arc has uncountably many composants. (Hint: each composant is a countable
union of proper subcontinua, each of which is nowhere dense by Theorem~\ref{thm:thin}; use the
Baire category theorem, Lemma~\ref{lem:baire}.)
\end{exercise}

\begin{exercise}
Using Corollary~\ref{cor:sub-homeo}, show that a pseudo-arc $M$ contains a strictly decreasing
sequence $M=K_0\supsetneq K_1\supsetneq K_2\supsetneq\cdots$ of subcontinua, each homeomorphic to
$M$, whose intersection is a single point.
\end{exercise}
